\clearpage\section*{混合自测：先遮住后面的解析}\addcontentsline{toc}{section}{混合自测题目（Q52--Q56）}
这五题均为自拟复习题，不表示考试比例或时限。先只看cheatsheet独立完成，圈出关键词并写下查表编号；从下一组Q52解析开始再核对。每题同时考查定位、计算与解释。

\subsection*{Q52 同一个输入，三种更新为什么不同}
查表：正面第2栏P01；第3栏A01、B01；第4栏B02。

单样本$x=1$、目标$z=1$，初始$w=.2,b=.1,\eta=.1$。分别把该单元解释为：(A) $\hat y=H(wx+b),H(0)=1$的感知机；(B) $a=wx+b$的ADALINE，$E=\frac12(z-a)^2$；(C) $a=\sigma(wx+b)$的单个输出神经元，仍用半平方误差。求三者一次同步更新后的$w,b$，说明为何仅凭“更新权重”不能选公式。

\subsection*{Q53 有核系数就一定是最优分类器吗}
查表：背面第1栏S01、S02；第2栏K01。

线性核，训练$x=(-2,2)$、$y=(-1,1)$。有人交出$\alpha_1=\alpha_2=1/8,b=1$，并称两点都是自由支持向量。检查可行性，纠正$b$，再求$w$、总间隔和primal/dual值。若将$(w,b)$再乘2，哪些量保持不变？这里按hard-margin讨论。

\subsection*{Q54 加对角能否把无效核变成原来的有效核}
查表：背面第2栏K01；第4栏L02；第5栏L03。

拟用$K=\left[\begin{smallmatrix}1&2\\2&1\end{smallmatrix}\right]$作为两点Gram矩阵，$y=(0,2)^T,\gamma=1/2$。先找一个向量证明该Gram无效，再检查$H=K+2I$是否正定。若形式上求带bias系统，会得到什么解？能否据此宣布原$K$是有效核？

\subsection*{Q55 改变管宽后，原来的KKT证据还成立吗}
查表：背面第3栏V01、V02。

一条样本$x=0,y=3$，线性核，$C=1$。先取$\varepsilon=0$，候选$w=0,b=3,\xi=\xi^*=0,\alpha=\alpha^*=2/5$。逐项检查KKT。再改$\varepsilon=1/2$而保持这些变量不变：原预测是否仍最优？原乘子是否仍是有效KKT证据？给出可行的修正乘子。

\subsection*{Q56 训练目标最优，能否直接宣布泛化最好}
查表：背面第4栏L01；第5栏L03；正面第4栏B03。

训练样本$(0,0),(1,2)$，线性LS-SVM，$\gamma=1$，bias不惩罚。候选A为$f_A(x)=2x$，候选B为$f_B(x)=2x/3+2/3$。计算两者训练目标。另有一个此前未用来拟合的测试样本$(2,4)$，比较该点平方预测误差。解释两种排序为什么可相反，以及这个单点能支持多强的泛化结论。

\problem{Q52}{混合自测解析：同分数、不同学习信号}{正面第2栏P01；第3栏A01、B01；第4栏B02}
\qtext{数据$x=z=1,w=.2,b=.1,\eta=.1$。对硬阈值感知机、线性ADALINE、sigmoid输出加半平方误差三种设定，各作一步同步更新。初始净输入均为$u=.3$。}
\route{“硬类别”指P01；“线性连续输出的平方残差”指A01；“sigmoid加链式导数”指B02。三者有同一个$u$，但预测和损失对参数的依赖不同。}
\given{$w,b$均为标量；$H(u)=1$当$u\ge0$；$\sigma(u)=1/(1+e^{-u})$；每种情况独立从相同旧参数开始，不能顺接更新。}
\solve
(A) 感知机：$\hat y=H(.3)=1$，$z-\hat y=0$，所以$w^+=.2,b^+=.1$。

(B) ADALINE：$a=.3$，线性残差$e=.7$，所以$\Delta w=.1(.7)(1)=.07$、$\Delta b=.07$，得到$w^+=.27,b^+=.17$。

(C) sigmoid输出：$a=\sigma(.3)\approx.57444252$。使用课件负导数约定，
\[
\delta=(z-a)a(1-a)\approx.10403106,
\quad\Delta w=\eta\delta x\approx.01040311,\quad\Delta b\approx.01040311.
\]
故$w^+\approx.21040311,b^+\approx.11040311$。这里是一个输出单元，无需额外编造隐藏层。
\ans{三者新参数分别为$(.2,.1)$、$(.27,.17)$、约$(.21040311,.11040311)$。相同输入和学习率不能确定更新式；必须核对输出激活及训练误差信号。}
\qcheck{感知机分类已经正确，零更新合理；LMS更新后$a=.44$，半平方误差由$.245$降到$.1568$；sigmoid情形通过直接前向损失的有限差分可验证梯度符号。}
\pitfall{C不能直接套$e=.7$，也不能只用$z-a$而漏$a(1-a)$。A不能因$u=.3\ne1$就误以为分类错误。}
\transfer{改为线性输出时删激活导数；目标若改成别的损失须重新链式求导。若新增隐藏层，还要用同一旧输出权重反传。}

\problem{Q53}{混合自测解析：偏置、间隔与最优性一起检查}{背面第1栏S01、S02；背面第2栏K01}
\qtext{线性核，$x=(-2,2),y=(-1,1),\alpha=(1/8,1/8)$，候选$b=1$。检查并纠正hard-margin解，再分析正比例缩放。}
\route{圈“hard-margin、给乘子、检查偏置”。S02从系数恢复$w$并检查KKT，S01检查几何间隔；只看到乘子非负不足以证明最优。}
\given{$w=\sum\alpha_i y_ix_i$为标量；函数间隔$m_i=y_i(wx_i+b)$；无软间隔$C$，题中的“自由SV”应理解为hard中$\alpha_i>0$。}
\solve
先算$w=(1/8)(-1)(-2)+(1/8)(1)(2)=1/2$。若$b=1$，则$g=(0,2)$，函数间隔$(0,2)$，第一点违反$m\ge1$，因此不可行。

两点乘子均正，应满足$m=1$。用任一点恢复：
\[
b=y_1-wx_1=-1-(1/2)(-2)=0,
\quad y_2-wx_2=1-1=0.
\]
正确解为$w=1/2,b=0$，两间隔均1，总宽$2/|w|=4$。primal为$\frac12w^2=1/8$；dual为
\[
\sum\alpha_i-\frac12\left(\sum\alpha_i y_i x_i\right)^2
=\frac14-\frac18=\frac18.
\]
非负、$\sum\alpha_i y_i=0$、原可行、驻点、互补均满足，且原对偶同值。

把$(w,b)$乘2，变为$(1,0)$：决策边界仍$x=0$，分类不变，几何间隔仍2；函数间隔从1变2。此时数据不再处在规范的$g=\pm1$支撑面上，不能把$2/\|w\|=2$当本数据的最大总宽。
\ans{错误是$b=1$；正确$b=0,w=1/2$，总间隔4，原/对偶最优值均$1/8$。正比例缩放保留边界和几何距离，改变函数间隔与规范化。}
\qcheck{直接点到面距离$|x_i|=2$，两侧总宽4，与规范公式一致。}
\pitfall{“已给$\alpha$”不等于给定答案正确；必须连同$b$检查全部约束。}
\transfer{soft-margin求$b$要改用$0<\alpha<C$；核化后通常不显式求$w$，用$\sum\alpha_i y_iK_{is}$恢复$b$。}

\problem{Q54}{混合自测解析：能解方程不等于核合法}{背面第2栏K01；第4栏L02；第5栏L03}
\qtext{$K=\left[\begin{smallmatrix}1&2\\2&1\end{smallmatrix}\right]$、$\gamma=1/2,y=(0,2)^T$。证明$K$无效，检查$H=K+2I$并形式求解，然后指出不能得出的结论。}
\route{圈“Gram有效性、加对角、能否宣布”。先K01检查原核，再L02做代数；L03的唯一性充分条件不能反向推回原核必有效。}
\given{$c=(1,-1)^T$用作PSD反例；$H=\left[\begin{smallmatrix}3&2\\2&3\end{smallmatrix}\right]$，未知仍$[\alpha_1,\alpha_2,b]$。这里只作错误诊断，不将无效核纳为标准LS-SVM。}
\solve
$c^TKc=-2<0$，所以$K$不是PSD，不能是实内积特征的Gram矩阵。$H$的特征值是$3+2=5$和$3-2=1$，均正，故本矩阵$H$确实正定。

形式系统为
\[
\left[\begin{array}{ccc|c}3&2&1&0\\2&3&1&2\\1&1&0&0\end{array}\right].
\]
末行给$\alpha_1=-\alpha_2$；令$\alpha_2=a$，前两行化为$-a+b=0$、$a+b=2$，因此$a=1,b=1$，即$\alpha=(-1,1)^T$。

这说明一个给定代数系统可解；并没有改变原$K$的负二次型。若把$H$本身重新定义为另一个有限集合上的核矩阵，那已是换模型，也仍要说明完整核函数及适用域，不能声称原$K$被“证明有效”。
\ans{原$K$无效；$H$正定且形式系统解为$\alpha=(-1,1),b=1$。可逆/可解性质不能替代有效核要求。}
\qcheck{$K\alpha+b\mathbf1=(2,0)^T$，形式残差$(-2,2)^T=\alpha/\gamma$；代数条件成立，但$c^TKc=-2$仍未消失。}
\pitfall{不能从一个有限点集的修补矩阵PSD推出原核函数对所有输入有效；更不能从对称性单独推出PSD。}
\transfer{若换成有效$K\succeq0$且$\gamma>0$，L03的保证可直接使用。若Gram计算漏了特征映射的$\sqrt2$项，先重算正确核，不应任意改数字掩盖错误。}

\problem{Q55}{混合自测解析：预测最优与乘子证据不同}{背面第3栏V01、V02}
\qtext{一条样本$x=0,y=3$，线性SVR，$C=1$；候选$w=0,b=3,\xi=\xi^*=0,\alpha=\alpha^*=2/5$。先验$\varepsilon=0$，再换为$1/2$；判断原预测和原乘子是否均可保留。}
\route{圈“$\varepsilon=0$、两乘子同正、换管宽”。V02明示零管宽边界需重判；要区分原始变量的最优性和给定乘子的KKT有效性。}
\given{$r=y-f=0$，$d=\alpha-\alpha^*=0$，核$K(0,0)=0$；由驻点可取松弛非负约束的乘子$\eta=\eta^*=C-2/5=3/5$。}
\solve
当$\varepsilon=0$：两个误差约束都是$0\le0$，松弛非负；乘子在盒内；$\sum d=0$，$w=dx=0$。两条误差约束互补为$(2/5)\cdot0=0$，松弛互补同为0。原目标为0，对偶目标也为0，故KKT与最优性均成立。

改$\varepsilon=1/2$后，原始预测仍有$r=0$、零松弛、$w=0$，目标0已经是非负目标的下界，因此\textbf{预测仍最优}。然而误差约束有严格余量$1/2$，原乘子的互补乘积变为$(2/5)(1/2)=1/5\ne0$，所以\textbf{原乘子失效}。改为$\alpha=\alpha^*=0,\eta=\eta^*=1$即可满足全部条件。
\ans{$\varepsilon=0$允许同一点两乘子同时正；增到$1/2$时原预测仍最优，但必须改乘子，取两乘子0是一个正确选择。}
\qcheck{增管宽而保留原乘子会给dual值$-\varepsilon(\alpha+\alpha^*)=-2/5$，与primal值0不相等；改成零乘子后二者同为0。}
\pitfall{不能因为预测没有改变就跳过互补检查；也不能把“$\varepsilon>0$不能两乘子同正”无条件套在零管宽。}
\transfer{若残差不为0，增管宽还可能改变原始最优解；重新算管外损失、约束和KKT。无自由SV时使用可行bias区间，不随意套边界等号。}

\problem{Q56}{混合自测解析：目标与泛化评价可以排序相反}{背面第4栏L01；第5栏L03；正面第4栏B03}
\qtext{训练$(0,0),(1,2)$、线性LS-SVM、$\gamma=1$。比较A:$f=2x$与B:$f=2x/3+2/3$的训练目标；再比较未参与拟合的测试点$(2,4)$的平方误差，并解释结论边界。}
\route{圈“范数加平方误差、测试样本”。L01决定训练目标，B03决定如何解释测试结果。单看训练残差为0不能宣布A的LS目标更优。}
\given{A有$w_A=2,b_A=0$；B有$w_B=b_B=2/3$。训练目标$J=\frac12w^2+\frac12\sum e_i^2$；测试指标为$(y-f)^2$，题目未在此加$1/2$。}
\solve
A在两个训练点完全插值，$e=(0,0)$，但范数代价$\frac12(2)^2=2$，所以$J_A=2$。

B的训练残差$(-2/3,2/3)$，所以
\[
J_B=\frac12\frac49+\frac12\left(\frac49+\frac49\right)=\frac23<2.
\]
B是Q47已由最优条件验证的该训练问题解。对新点$x=2$，$f_A(2)=4$，平方误差0；$f_B(2)=2$，平方误差$(4-2)^2=4$。于是训练正则目标偏好B，但给定测试点偏好A；没有矛盾，因为评价对象和依据不同。

一个测试点只能报告这个点的结果，不能证明A对总体分布一定更好。模型和超参数应使用训练/验证过程选择；测试数据保留最终评价，不能在观察结果后反复调参还把它当独立测试。
\ans{$J_A=2,J_B=2/3$；测试平方误差分别0和4。训练目标最优不保证每个未见点更准，单点测试也不足以确定整体泛化优劣。}
\qcheck{B的残差与$\alpha/\gamma$一致；A虽然训练误差零，但$w=2$且若$e=0$则$\alpha=0$，不满足$w=\sum\alpha_i x_i$，所以不是该正则问题最优解。}
\pitfall{别把“平方误差”默认都带$1/2$；本题明确区分训练目标与测试指标。也不要把用户提供的“最大”误读成要最大化这些损失。}
\transfer{换$\gamma$会改变训练权衡；换测试分布会改变泛化排序。保持“先定义指标、再计算、再限定结论”的步骤，不能仅凭模型名称比较好坏。}
